% !TeX program = xelatex
% Слайды к докладу на НСКФ'2026 на основе svf_so_nscf.tex
% Компиляция:  xelatex svf_so_slides.tex   (второй проход — для номеров слайдов)

\documentclass[10pt,aspectratio=169]{beamer}

\usetheme{metropolis}

% \makeatletter
% \def\beamer@framenotesbegin{%
%     \usebeamercolor[fg]{normal text}%
%     \gdef\beamer@noteitems{}%
%     \gdef\beamer@notes{}%
% }
% \makeatother

\usepackage[utf8]{inputenc}
\usepackage[russian,english]{babel}

\usepackage{amsmath,amssymb,bm}
\usepackage{booktabs}
\usepackage{xcolor}
% \usepackage{pdfcomment}
% \setbeameroption{show notes}
% \let\oldnote\note
% \renewcommand{\note}[1]{%
%   \pdfcomment[author=Speaker, color=0 0 0, opacity=1]{#1}%
% }
% ——— Обозначения по тексту статьи svf_so_nscf.tex ———
\usepackage{xspace}

\newcommand{\LRb}[1]{\left(#1\right)}
\newcommand{\LRB}[1]{\Bigl(#1\Bigr)}
\newcommand{\LRC}[1]{\Bigl\{#1\Bigr\}}
\newcommand{\SET}[1]{\left\{#1\right\}}
\newcommand{\FRAC}[2]{\frac{\displaystyle #1}{\displaystyle #2}}
\newcommand{\MOD}[1]{\left|#1\right|}
\newcommand{\abs}[1]{\left|#1\right|}
\newcommand{\norm}[1]{\left\|#1\right\|}
\newcommand{\LE}{{\leqslant}}
\newcommand{\GE}{{\geqslant}}
\newcommand{\EQ}{{=}}
\newcommand{\IN}{{\in}}
\newcommand{\RR}{\mathbb{R}}
\newcommand{\eps}{\varepsilon}
\newcommand{\EI}{K}
\newcommand{\LK}{L}
\newcommand{\CV}{\mathrm{CV}}
\newcommand{\MSD}{\mathrm{MSE}}
\newcommand{\FS}{\mathcal{R}}
\newcommand{\argmin}{\mathop{\mathrm{arg\,min}}}
\newcommand{\argmax}{\mathop{\mathrm{arg\,max}}}
\newcommand{\dima}{m}%{{n_{\alpha}}} % dimension of alpha

\newcommand{\spotoptim}{\texttt{spotoptim}\xspace}
\newcommand{\oldoptim}{\texttt{oldoptim}\xspace}

% ——— Нижний колонтитул ———
\newcommand{\conferenceFooter}{%
  НСКФ'2026. Суррогатная оптимизация в SvF-технологии, %
  Волошинов В.В., Куликов Ф.Е., Соколов А.В.}

\setbeamersize{text margin left=0.35cm, text margin right=0.35cm}
\setbeamerfont{footline}{size=\fontsize{6}{7}\selectfont}
\setbeamercolor{footline}{fg=structure.fg!60!black, bg=white}

\setbeamertemplate{footline}{%
  \leavevmode%
  \begin{beamercolorbox}[wd=\paperwidth,ht=2.1ex,dp=0.5ex,%
    leftskip=0.35cm,rightskip=0.35cm plus1fil]{footline}%
    \usebeamerfont{footline}\conferenceFooter\hfill\insertframenumber/\inserttotalframenumber%
  \end{beamercolorbox}%
}

% \setbeameroption{show notes on second screen=right} 
% \setbeamertemplate{note page}{%
%   \insertnote
% }

% ——— Титульный слайд ———
\title[уррогатная оптимизация в SvF-технологии]{%
  Сравнение алгоритмов суррогатной оптимизации\\ в SvF-технологии сбалансированной идентификации \\
  \textcolor{gray}{\small Докладчик: Куликов Филипп Евгеньевич} \\[-1.2ex]
  \textcolor{gray}{\small Авторы:  В. В. Волошинов, Ф. Е. Куликов, А. В. Соколов} \\[-1.2ex]
  \textcolor{gray}{\small Организация:  Центр распределённых вычислений ИППИ РАН}}
\author[Волошинов, Куликов, Соколов]{%
  В. В. Волошинов, Ф. Е. Куликов, А. В. Соколов}
\institute[ИППИ РАН]{Институт проблем передачи информации им.\ А.\,А.~Харкевича РАН, Москва}
\date{НСКФ'2026}

\begin{document}

% ============================ 1 ============================
\begin{frame}[plain]
  \usebeamerfont{title}\usebeamercolor[fg]{title}\inserttitle
\end{frame}

% ============================ 2 ============================
\begin{frame}{Структура доклада}
    SvF-технология (Simplicity vs Fitting) развивается больше 10 лет,
      применяется в прикладных задачах идентификации математических моделей.

    План презентации:

    \begin{itemize}
    \item Описание принципов SvF технологии
    \item Базовый оптимизационный алгоритм
    \item Теоретические основы имплементированного алгоритма
    \item Сравнение эффективности алгоритмов
    \item Выводы
  \end{itemize}

\note[item]{Emphasize the first bullet point.}
  \note[item]{Mention that handouts omit speaker notes.}
\end{frame}

% ============================ 3 ============================
\begin{frame}{Описание SvF}
  SvF технология позволяет получить интерпретируемую модель для описания набора данных.

  Пример модели в виде системы ОДУ второго порядка для линейного маятника:

  \begin{columns}[T]
    \begin{column}{0.5\textwidth}
      \vspace{1cm}
      $$
      \begin{cases}
      \ddot{x} = f(x, v) \\
      \dot{x} = v
      \end{cases}
      $$

      $$
      F(x, K, \alpha) = \frac{1}{|K|} \sum_{k \in K} (z_k - x(t_k))^2
       + \alpha \int_T \left( \frac{d^2 x}{dt^2} \right)^2 dt
      $$
    \end{column}
    \begin{column}{0.4\textwidth}
      \begin{center}
        \includegraphics[width=\textwidth]{images/Regionx.png}\\[-0.5ex]
      \end{center}
    \end{column}
  \end{columns}

%  \[
%    Q \EQ \SET{x\in [-0.1, 2.2], v \in [-1, 1.5] : f_i(x)\LE 0\ (i\in I),\ g_j(x)\EQ 0\ (j\in J)}.
%  \]

\note {ЗАМЕТКИ ДЛЯ СЛАЙДА}
\end{frame}


% ============================ 4 ============================

\begin{frame}{Нижний уровень оптимизации в SvF}

Типичным критерием выбора модели является минимизация функционала
$$
F(x,\EI^\prime,\alpha) = F_D(x,\EI^\prime) + \FS(\alpha ,x) \to \min\limits_{x\in Q},
$$

$$
\FS(\alpha ,x) =  \iint [ \alpha_x^2 \left( \frac{\partial^2 f}{\partial x^2} \right)^2
+ 2\alpha_x \alpha_v \left( \frac{\partial^2 f}{\partial x \partial v} \right)^2 + \alpha_v^2 \left( \frac{\partial^2 f}{\partial v^2} \right)^2 ] dx dv.
$$


  Разбиение $\EI = E \cup T$, $E\cap T=\varnothing$. Для заданного $\alpha$:
  \[
    x_E^\alpha = \argmin_{x\in Q} F(x,E,\alpha),
    \qquad
    \CV^\alpha_{E,T} = \MSD(x_E^\alpha, T).
  \]
  По $\LK$ парам $(E_\ell,T_\ell)$ --- агрегированная CV-ошибка 
  $\CV(\alpha)=\FRAC{1}{\sum_\ell \abs{T_\ell}}\sum_\ell \abs{T_\ell}\CV^\alpha_{E_\ell,T_\ell}$.

  \begin{itemize}\setlength{\itemsep}{1pt}
    \item Вызовы решения целого набор задач дорогостоящие
    \item Задачи нижнего уровня независимы --- параллелятся.
  \end{itemize}
\end{frame}


\begin{frame}{Предсказания модели при различных $\alpha$}
  \begin{center}
    \includegraphics[width=0.95\textwidth]{images/comparison_of_apha_2.png}\\[-0.5ex]
  \end{center}
\end{frame}

% ============================ 5 ============================
\begin{frame}{Базовый алгоритм: квадратичный суррогат}
  По накопленным парам $\SET{\LRB{\CV(\alpha^\nu),\,\alpha^\nu}}_{\nu\EQ 1}^N$
  строится многочлен 2-й степени
  \[
    \Pi(\bm{\pi},\alpha)=\sum_{1\LE i\LE j\LE m}\pi_{ij}\alpha_i\alpha_j
      +\sum_{1\LE i\LE m}\pi_i\alpha_i+\pi_0,
    \qquad \bm{\pi}\in\RR^{\frac{(m+2)(m+1)}{2}}.
  \]
  Коэффициенты --- из сильно выпуклой регуляризованной задачи
  \[
    \sum_{\nu\EQ 1}^{N}\FRAC{\LRb{\CV(\alpha^\nu)-\Pi(\bm{\pi},\alpha^\nu)}^2}
      {e^{w\norm{\alpha^\nu-\alpha^N}}}
    \;+\;\mu\LRB{\sum_{i\LE j}\pi_{ij}^2}\;\to\!\min_{\bm{\pi}} .
  \]
  Новая точка --- локальный минимум суррогата в пределах допустимого радиуса:
  \[
    \Pi(\bm{\pi}^*,\alpha)\to\min_{\alpha\in\RR_+^m,\ \norm{\alpha-\alpha^N}\LE s}.
  \]
  Погрешность $\Delta(w,\mu,s)=\MOD{\CV(\alpha^*)-\Pi(\bm{\pi}(w,\mu),\alpha^*)}$
  используется для настройки $w,\mu$.
\end{frame}

% ============================ 6 ============================
\begin{frame}{Базовый алгоритм: эвристики}
  \begin{itemize}
    \item Начальная последовательность --- малые возмущения одного
      заданного вектора $\alpha$ ($m+1$ точка).
    \item Параметры $w$, $\mu$ поочерёдно мультипликативно корректируются
      ($\mu'=(1+d)\mu$, $d\EQ 0.1\%$).
    \item Радиус доверия $s$: сначала «агрессивно» растёт,
      далее --- по эмпирическим правилам.
    \item \textbf{Проблема «оврага»}: при $m>1$ точки прыгают со «склона»
      на «склон». Признак --- большой косинус между соседними
      $\Delta^\nu\alpha=\alpha^\nu-\alpha^{\nu-1}$.
      Реакция --- шаг вдоль усреднённого направления.
  \end{itemize}
  \vspace{0.5em}

   \textbf{Реализация:} \texttt{https://github.com/distcomp/SvF/SvFlib/SurMin.py}.
\end{frame}

% ============================ 7 ============================
\begin{frame}{Описание алгоритма в spotoptim}
  $Y(x)$ --- реализация гауссова процесса; корреляция значений:
  \[
    \mathrm{Corr}\bigl[Y(x_i),Y(x_j)\bigr]
      = \exp\Bigl(-\sum_{\ell=1}^{m}\theta_\ell \abs{x_{i\ell}-x_{j\ell}}^{p_\ell}\Bigr),
    \qquad p_\ell\EQ 2.
  \]

$$
L({Y(x_1)}, \ldots, {Y(x_n)} | \mu, \sigma) = \frac{1}{(2\pi\sigma^2)^{n/2} |\mathbf{R}|^{1/2}} \exp \left[ -\frac{(\mathbf{y} - \mathbf{1}\mu)^T \mathbf{R}^{-1} (\mathbf{y} - \mathbf{1}\mu)}{2\sigma^2} \right],
$$

  % Параметры $\mu$, $\sigma^2$, $\bm{\theta}$ --- по максимальному правдоподобию;
  % после исключения $\hat\mu$, $\hat\sigma^2$ максимизируется
  % \[
  %   -\frac{n}{2}\log\hat\sigma^2 - \frac12 \log\abs{\mathbf{R}}
  %   \;\to\max_{\bm{\theta}}
  % \qquad\text{(генетический алгоритм).}
  % \]
  Предсказание в новой точке $x^*$:
  $$
    \hat y(x^*) = \hat\mu + \mathbf{r}^T\mathbf{R}^{-1}(\mathbf{y}-\mathbf{1}\hat\mu),
  $$
  $$
\hat{y}(\mathbf{x}^{*}) = a + \sum_{i=1}^{n} b_{i} \, \varphi(\mathbf{x}^{*} - \mathbf{x}_{i})
$$
где $\phi(z)$ - радиальная базисная функция вида
$$
\varphi(\mathbf{z}) = \exp\left(-\sum_{\ell=1}^{\dima} \theta_{\ell} \, |z_{\ell}|^{p_{\ell}}\right) 
$$
\end{frame}
% ============================ 8 ============================
\begin{frame}{Функции приобретения}
  Точка нового запуска $\CV(\alpha)$ выбирается по суррогатной модели.
  В \spotoptim три варианта, тестировались два:

  \vspace{0.6em}
  \begin{description}\setlength{\itemsep}{0.6em}
    \item[Математическое ожидание (\spotoptim-y)]
      $\alpha^{*}=\argmin \hat y(\alpha)$, дисперсия игнорируется.
      Быстрая сходимость, глобальный минимум не гарантируется.
      Глобальный поиск --- дифференциальная эволюция.

    \item[Expected Improvement (\spotoptim-ei)]
    $$
    EI(\mathbf{x}) = \int_{I=0}^{I=\infty} I \left\{ \frac{1}{\sqrt{2\pi s(\mathbf{x})}} \exp \left[ -\frac{(f_{\text{min}} - I - \hat{y}(\mathbf{x}))^2}{2s^2(\mathbf{x})} \right] \right\} dI.
    $$
  \end{description}
  \vfill
  Базовый алгоритм: детерминированный, без оценки неопределённости.
\end{frame}

\begin{frame}{Визуализация модели в spotoptim, mode=y}
  \begin{center}
    \includegraphics[width=0.6\textwidth]{images/figure_acquisition_mode_y.png}\\[-0.5ex]
  \end{center}
\end{frame}

\begin{frame}{Визуализация модели в spotoptim, mode=ei}
  \begin{center}
    \includegraphics[width=0.6\textwidth]{images/figure_acquisition_mode_ei.png}\\[-0.5ex]
  \end{center}
\end{frame}


% ============================ 9 ============================
\begin{frame}{Методика сравнения}
  Здесь и далее базовый алгоритм помечен как oldoptim

  Сравниваются три метода:
  \[
    \oldoptim,\qquad \spotoptim\text{-}y,\qquad \spotoptim\text{-}ei.
  \]
  \begin{enumerate}\setlength{\itemsep}{2pt}
    \item В каждой задаче --- фиксированный лимит итераций (60).
    \item \spotoptim: равномерный набор начальных точек в прямоугольнике
      допустимых значений $\alpha$.
    \item \oldoptim: одна стартовая точка (центр области).
    \item \textbf{Контроль:} \oldoptim перезапускается из лучшей точки
      начального набора \spotoptim.
  \end{enumerate}
  \vspace{0.4em}

\end{frame}

% ============================ 10 ============================
\begin{frame}{Тест 1. Линейный маятник, 3 коэффициента регуляризации}
  \[
    \frac{d^2x}{dt^2}=f(x,v),\quad v=\frac{dx}{dt},
    \qquad f\text{ --- полином 6-й степени от } (x,v).
  \]
  Итераций: 60; начальных точек \spotoptim: 12, область $[0.01,1]$;
  старт \oldoptim: $[0.5,0.5,0.5]$.
  \begin{center}
    \includegraphics[width=0.72\textwidth]{images/Osc-1.pdf}\\[-0.5ex]
    {\footnotesize Лучшее значение $\CV(\alpha)$ по номеру итерации.}
  \end{center}
\end{frame}

% ============================ 11 ============================
\begin{frame}{Тест 1: результаты}
  % \begin{table}\centering\small
  %   \begin{tabular}{lcc}
  %     \toprule
  %     \textbf{Метод} & $\CV$ \textbf{ошибка} & \textbf{Веса $\alpha$}\\
  %     \midrule
  %     \oldoptim & 46.78 & $[0.292,\,0.519,\,0.515]$\\
  %     \spotoptim-y & \textbf{17.98} & $[0.113,\,0.996,\,0.619]$\\
  %     \spotoptim-ei & 17.98 & $[0.116,\,0.865,\,0.946]$\\
  %     \oldoptim\ (хороший старт) & 17.99 & $[0.119,\,0.944,\,0.583]$\\
  %     \bottomrule
  %   \end{tabular}
  % \end{table}
  \begin{center}
    \includegraphics[width=0.6\textwidth]{images/Osc-2.pdf}\\[-0.5ex]
  \end{center}
  \vspace{0.6em}
  \begin{itemize}\setlength{\itemsep}{2pt}
    \item Из «центра» базовый алгоритм сходился медленно и
      остановился втрое хуже ($46.8$ против $18.0$).
    \item С удачной начальной точкой --- результат практически тот же,
      что у обоих \spotoptim.
    \item Две стратегии \spotoptim дали почти идентичный итог.
  \end{itemize}
\end{frame}

% ============================ 12 ============================
\begin{frame}{Тест 2. Транспирация сосны, 2 коэффициента регуляризации}
  Аппроксимируется сеточная функция $E(Q,VPD)$ --- транспирация от
  солнечной радиации $Q$ и дефицита давления пара $VPD$.
  Дифференциальных уравнений нет.
  Итераций: 60; начальных точек \spotoptim: 10, область $[0,1]$;
  старт \oldoptim: $[0.5,0.5]$.
  \begin{center}
    \includegraphics[width=0.72\textwidth]{images/E-1.pdf}\\[-0.5ex]
    {\footnotesize Лучшее значение $\CV(\alpha)$ по номеру итерации.}
  \end{center}
\end{frame}

% ============================ 13 ============================
\begin{frame}{Тест 2: результаты}
  % \begin{table}\centering\small
  %   \begin{tabular}{lcc}
  %     \toprule
  %     \textbf{Метод} & $\CV$ \textbf{ошибка} & \textbf{Веса $\alpha$}\\
  %     \midrule
  %     \oldoptim & 48.53 & $[0.112,\,0.186]$\\
  %     \spotoptim-y & \textbf{48.52} & $[0.099,\,0.187]$\\
  %     \spotoptim-ei & 48.55 & $[0.118,\,0.186]$\\
  %     \oldoptim\ (хороший старт) & 48.52 & $[0.097,\,0.189]$\\
  %     \bottomrule
  %   \end{tabular}
  % \end{table}
    \begin{center}
    \includegraphics[width=0.6\textwidth]{images/E-2.pdf}\\[-0.5ex]
  \end{center}
  \vspace{0.6em}
  \begin{itemize}\setlength{\itemsep}{2pt}
    \item Различия --- в четвёртом знаке: все методы нашли один и тот же оптимум.
    \item Базовый алгоритм сходился медленнее по итерациям,
      но монотонно и без «оврага».
    \item При $m=2$ (без дифференциальных уравнений) базовый алгоритм
      полностью работоспособен.
  \end{itemize}
\end{frame}

% ============================ 14 ============================
\begin{frame}{Тест 3. Три планеты, 10 коэффициентов регуляризации}
  Идентифицируются массы и закон притяжения $f(r)$ по зашумлённым
  траекториям трёх тел ($\approx 6\%$ гауссова шума):
  \[
    \textstyle
    \ddot r_i(t)=G\sum_{j\neq i}\frac{r_j(t)-r_i(t)}{\norm{r_i(t)-r_j(t)}^3},
    \qquad i=1{:}3 .
  \]
  $f(r)$ --- полином 8-й степени, $f(1)=1$.
  Итераций: 60; начальных точек \spotoptim: 20, область $[0,0.1]$.
  \begin{center}
    \includegraphics[width=0.66\textwidth]{images/3Body-1.pdf}\\[-0.5ex]
    {\footnotesize Базовый алгоритм остановился преждевременно.}
  \end{center}
\end{frame}

% ============================ 15 ============================
\begin{frame}{Тест 3: результаты}
  \begin{table}\centering\small
    \begin{tabular}{lc}
      \toprule
      \textbf{Метод} & $\CV$ \textbf{ошибка}\\
      \midrule
      \oldoptim & 12.64\\
      \spotoptim-y & 4.12\\
      \spotoptim-ei & \textbf{4.12}\\
      \oldoptim\ (хороший старт, 32 точки) & 4.12\\
      \bottomrule
    \end{tabular}
  \end{table}
  \vspace{0.6em}
  \begin{itemize}\setlength{\itemsep}{2pt}
    \item Размерность $m=10$ --- базовый алгоритм «застрял»: ошибка втрое выше.
    \item Даже с 32 начальными точками улучшение составило лишь
      два десятка итераций, затем остановка при
      $\norm{\Delta\alpha}<10^{-7}$.
    \item \spotoptim уверенно работает в высокой размерности ---
      главное практическое преимущество.
  \end{itemize}
\end{frame}

% ============================ 16 ============================
\begin{frame}{Выводы из сравнения}
  \begin{table}\centering
    \begin{tabular}{lccc}
      \toprule
      & $m\EQ 2$ & $m\EQ 3$ & $m\EQ 10$\\
      \midrule
      \oldoptim\ (центр области) & \checkmark & $\times$ & $\times$\\
      \oldoptim\ (хороший старт) & \checkmark & \checkmark & \checkmark\\
      \spotoptim-y, \spotoptim-ei & \checkmark & \checkmark & \checkmark\\
      \bottomrule
    \end{tabular}
  \end{table}
  \vspace{0.5em}
  \begin{itemize}\setlength{\itemsep}{2pt}
    \item При малом числе весов базовый алгоритм \textbf{не хуже} \spotoptim
    \item Ключевой фактор для базового алгоритма --- \textbf{выбор начальной точки},
      ранее явно не изучавшийся.
    \item Равномерное покрытие области начальных точек удобно и доступно пока только 
    вместе с spotoptim.
    \item \spotoptim лучше показывает себя при росте количества 
    коэффициентов регуляризации.
  \end{itemize}
\end{frame}

% ============================ 17 ============================
\begin{frame}{Результаты}
  \begin{columns}[T]
    \begin{column}{0.48\textwidth}
      \textbf{Подтверждено}
      \begin{itemize}\setlength{\itemsep}{1pt}
        \item Базовый алгоритм работоспособен.
        \item Результаты предыдущих публикаций, полученные с ним,
          корректны.
      \end{itemize}
    \end{column}
    \begin{column}{0.48\textwidth}
      \textbf{Расширено}
      \begin{itemize}\setlength{\itemsep}{1pt}
        \item Выбор двух дополнительных алгоритмов суррогатной оптимизации из spotoptim.
        \item Блок \texttt{COMPARE} в .mng --- графики сравнения методов.
        \item Несколько расчётов с разными стартами из RES-файлов
          в одном задании.
      \end{itemize}
    \end{column}
  \end{columns}
\end{frame}

% ============================ 18 ============================
\begin{frame}{}
  \centering\Large Спасибо за внимание!
\end{frame}

\end{document}
